% =====================================================================
%  The ghost-capture gap in declarative skeleton consistency.
%
%  BUILD (from mechanization/docs/):
%      TEXINPUTS=../../paper: pdflatex ghost-capture.tex     (twice)
%  or  make -f Makefile.ghost-capture
%
%  Standalone document: it has its own preamble and body, and reuses the
%  paper's notation by \input-ing paper/macro.tex, so every judgement and
%  every rule name below is typeset by the paper's own macros.
% =====================================================================
\documentclass[11pt,a4paper]{article}

\usepackage[utf8]{inputenc}
\usepackage[T1]{fontenc}
\usepackage{amsmath,amssymb,mathtools,amsthm}
\usepackage{graphicx}
\usepackage{thmtools}
\usepackage{mathpartir}
\usepackage{stmaryrd}
\usepackage{xspace}
\usepackage{url}
\usepackage{xcolor}
\usepackage{array}
\usepackage{tikz}
\usepackage[normalem]{ulem}
\usepackage{enumerate}
\usepackage{enumitem}
\usepackage{booktabs}
\usepackage{halloweenmath}
\usepackage{trimclip}
\usepackage{scalerel}[2016/12/29]
\usepackage[old]{old-arrows}
\usepackage[ruled, noend]{algorithm2e}
\usepackage{listings}
\usepackage{pifont}
\usepackage{fontawesome}
\usepackage[margin=1in]{geometry}
\usepackage{hyperref}
\usepackage{cleveref}

\input{macro}

\theoremstyle{definition}
\newtheorem{factc}{Fact}
\crefname{factc}{Fact}{Facts}
\Crefname{factc}{Fact}{Facts}

\newcommand{\coq}[1]{\texttt{\small #1}}
\newcommand{\defeq}{\mathrel{{:}{=}}}
\newcommand{\MC}{\textbf{[machine-checked]}\xspace}
\newcommand{\ARG}{\textbf{[argued]}\xspace}

% the configuration, written once
\newcommand{\Scex}{\forall\alpha.\Hole}
\newcommand{\Pcex}{\forall\alpha.\alpha}

\title{The ghost-capture configuration\\
  \large \Calc{} presupposes that skeleton consistency implies refinement.
  That presupposition is load-bearing, it is not derivable from the rules,
  and the mechanisation cannot go on deferring it}
\author{}
\date{}

\begin{document}
\maketitle

\begin{abstract}
  \noindent
\Calc{} declares, in prose, that skeleton consistency presupposes
  refinement --- ``\emph{skeleton consistency
  $\typconsis[\Psi]{\Gamma}{\syn{S}}{\bid{P}}$ requires well-formedness
  of $\syn{S}$ and $\bid{P}$ under $\Gamma$ \textbf{as well as the
  refinement judgement} $\Gamma \vdash \erase{S}\fill P$}''
  (\coq{declarative.tex:2074}).  \textbf{This note is not a criticism of
  that statement; it is a demonstration that it cannot be treated as
  bookkeeping.}  The mechanisation deferred it --- its own registry
  records the deferral as ``\emph{it has not bitten yet}''
  (\coq{docs/WF-SIDE-CONDITIONS.md} \S2) --- and it has now bitten, at
  the one place the two inference theorems have to cross.

  The presupposition is not implied by the rules.  The consistency rules push a
  quantifier's bound variable into $\Gamma$; the refinement rules push it
  into $\Delta$ and then forbid a ghost's fill to mention $\Delta$.  So
  the rules alone permit exactly the fill the side condition exists to
  forbid: a ghost \emph{captured} by a quantifier of its own subject.
  One configuration, $\syn{\Scex}$ against $\syn{\Pcex}$, exhibits it.

  Three consequences follow, each with a Coq witness.  Read without the
  presupposition, declarative skeleton inference derives
  $\pati{\Gamma}{x}{\syn{\Pcex}}$ at $\Gamma = x:\Scex$, which no run can
  produce.  The two inference theorems therefore cannot be restated over
  the declarative judgements so read.  And the obvious weakening ---
  exploiting the existential skeleton of \patlab{App} --- does not
  survive either.

  The repair is to put the paper's own sentence on the rules; no change
  to any published statement is proposed.  \Cref{sec:repair} gives two
  renderings and says what each costs, and \Cref{sec:survives} says which
  of the facts below the repair retracts --- \emph{that they retract is
  the point}.

  Everything marked \MC is \texttt{Qed} and axiom-free
  (\texttt{Print Assumptions}: \emph{closed under the global context}).
\end{abstract}

\begin{center}\fbox{\parbox{0.93\textwidth}{\small
\textbf{Status (2026-08-11).}  The repair of \Cref{sec:repair} has been
\textbf{applied}, in the form (R-a): the two presuppositions are now premises
of the declarative rules.  Everything in
\Cref{sec:gap,sec:inference,sec:app} describes the judgements read
\emph{without} them, and under that reading it stands; with them,
\Cref{fact:pat,fact:bridge,fact:app} are \textbf{retracted}, machine-checked at
the very configuration that produced them
(\coq{cap\_d\_pat\_excluded}, \coq{cap\_bridge\_holds}).
\Cref{fact:consis,fact:refine} are untouched --- see \Cref{sec:survives}.

\smallskip
What it bought: \coq{d\_pat\_to\_ad} and \coq{d\_sd\_typing\_to\_ad} are
theorems with \emph{no} hypothesis, and the closed-program completeness
corollaries are now stated over \Calc's own declarative judgements
($\sdtypp[\M]{\cdot}{M}{\bid{B}}$ and $\typp[\M]{\cdot}{M}{\bid{B}}$), matching
the soundness ones.  \coq{make coq} exit 0; 28/28 \texttt{Print Assumptions}
\emph{closed under the global context}; no \texttt{Admitted}, no
\texttt{Axiom}.  Cost: two premises, one new judgement
($\skctx[\Delta]{\Gamma} P \fill Q$, the uncoloured refinement the paper
already defines at \coq{declarative.tex:1516}) and its three structural
lemmas.  \emph{No published statement was changed.}}}\end{center}

% =====================================================================
\section{Summary}\label{sec:summary}
% =====================================================================

\begin{center}\small
\begin{tabular}{@{}p{6.1cm}p{3.0cm}p{5.2cm}@{}}
  \toprule
  Statement & Status & Coq witness \\
  \midrule
  $\typconsis[\cdot]{\Gamma}{\syn{\Scex}}{\syn{\Pcex}}$ is derivable
    & \MC
    & \coq{cap\_consis} \\[.4ex]
  $\skctx[\cdot]{\Gamma}\Scex \fill \Pcex$ is \emph{not}
    & \MC
    & \coq{cap\_no\_refine} \\[.9ex]
  \midrule
  \emph{hence} the paper's side condition is not derivable from the
  consistency rules
    & \MC
    & the two above \\[.9ex]
  \midrule
  $\pati{x:\Scex}{x}{\syn{\Pcex}}$ is derivable declaratively
    & \MC
    & \coq{cap\_d\_pat} \\[.4ex]
  \dots and refuted once the side condition is carried
    & \MC
    & \coq{cap\_not\_reachable} \\[.4ex]
  \emph{hence} declarative skeleton inference does not bridge to the
  side-condition-carrying one
    & \MC
    & \coq{d\_pat\_bridges\_false} \\[.9ex]
  \midrule
  \dots and \patlab{App}'s \emph{existentially quantified} skeleton does
  not repair the bridge
    & \MC
    & \coq{d\_pat\_app\_package\_false} \\[.9ex]
  \midrule
  the configuration satisfies both conditions of
  Def.~(problems and solutions of subtyping)
    & \ARG (\Cref{sec:solutions})
    & --- \\
  \bottomrule
\end{tabular}
\end{center}

\noindent
Files: \coq{cex/cex\_ad\_pat\_join\_not\_leaf\_valid.v} (the
configuration) and \coq{cex/cex\_d\_pat\_bridges\_false.v} (the two
bridge refutations).  \Cref{sec:survives} says which rows of this table
the repair of \Cref{sec:repair} retracts.

% =====================================================================
\section{What the paper says}\label{sec:paper}
% =====================================================================

Four statements, quoted or transcribed, are all this note needs.

\paragraph{(P1) Refinement tracks the subject's own binders, and forbids
ghosts to use them.}
\coq{declarative.tex:1422--1452}:

\begin{quote}\itshape
  We define a skeleton refinement relation
  $\Gamma \vdash \inh{P}\fillsyn\bid{Q}$ which fills in ghosts on the
  LHS with synthesised skeletons to obtain the RHS.\ \dots\ given a
  skeleton $P$, for every ghost in $P$, \textbf{it cannot represent
  skeletons that use type variables bound in $P$}.\ \dots\ To enforce
  this requirement, we augment the full judgement form to
  $\skctx[\Delta]{\Gamma}\inh{P}\fillsyn\bid{Q}$, where the set $\Delta$
  tracks type variables bound in $P$.
\end{quote}

\noindent
The two rules that carry the discipline are

\begin{mathpar}
\inferrule*[Lab=\text{\scshape R-Forall}]
{
  \skctx[\Delta,\alpha]{\Gamma,\alpha} \inh{P}\fillsyn \bid{Q}
}
{\skctx[\Delta]{\Gamma}
  \inh{\forall\alpha.P} \fillsyn \inh{\forall\alpha.}\bid{Q}}
\and
\inferrule*[Lab=\text{\scshape R-Ghost}]
{
  \Delta\cap\ftv{\syn{Q}} = \emptyset
}
{\skctx[\Delta]{\Gamma} \Hole \fillsyn \syn{Q}}
\end{mathpar}

\noindent
\emph{The bound variable goes into $\Delta$, and $\Delta$ is what a
ghost's fill must avoid.}

\paragraph{(P2) Consistency puts the same bound variable somewhere else.}
\coq{declarative.tex:1966--2045}.  The two rules this note needs are

\begin{mathpar}
\inferrule*[Lab=\eqlab{ForallSyn}]
{
  \typconsis[\Psi]{\Gamma,\alpha}{\syn{S}}{\bid{P}}
}
{\typconsis[\Psi]{\Gamma}
  {\syn{\foralltag{\delta}\alpha.S}}{\syn{\forall\alpha.} \bid{P}}}
\and
\inferrule*[Lab=\eqlab{Ghost}]
{~}
{\typconsis[\cdot]{\Gamma}{\Hole}{\inh{P}}}
\end{mathpar}

\noindent
\eqlab{ForallSyn} puts $\alpha$ into $\Gamma$; \eqlab{Ghost} matches a
ghost against \emph{any} inherited skeleton.  There is no $\Delta$, so
nothing stops the $\inh{P}$ of \eqlab{Ghost} from being the $\alpha$ that
\eqlab{ForallSyn} just bound.

\paragraph{(P3) The side condition that was meant to stop it.}
\coq{declarative.tex:2073--2075}:

\begin{quote}\itshape
  Skeleton consistency $\typconsis[\Psi]{\Gamma}{\syn{S}}{\bid{P}}$
  requires well-formedness of $\syn{S}$ and $\bid{P}$ under $\Gamma$
  \textbf{as well as the refinement judgement}
  $\Gamma \vdash \erase{S}\fill P$.
\end{quote}

\noindent
This is an \emph{external} requirement in prose, in the style the
mechanisation registers separately (\coq{docs/WF-SIDE-CONDITIONS.md}
\S2).  It is exactly the missing $\Delta$ discipline of (P1),
re-imposed from outside.  \Cref{sec:gap} shows it is not implied by the
rules.

\paragraph{(P4) Where consistency is consumed.}
The three base rules of declarative skeleton inference embed the
skeleton subtyping judgement, which bottoms out in consistency:

\begin{mathpar}
\inferrule*[Lab=\patlab{Var}]
{
  \Gamma\ni x:P \\
  \patsub[\M][\cdot]{\Gamma}{\syn{P}}{\bid{Q}}
}
{\patp[\M]{\Gamma}{x}{\bid{Q}}}
\and
\inferrule*[Lab=\sublab{Infer}]
{
  \typconsis[\cdot]{\Gamma}{\syn{S}}{\syn{P}}
}
{\patsub[\TI][\cdot]{\Gamma}{\syn{S}}{\syn{P}}}
\end{mathpar}

\noindent
(\coq{declarative.tex:1634--1653}, \coq{:1897}; \patlab{Unit} and
\patlab{Anno} are the same shape with $\tunit$ and with the
annotation's skeleton as subject.)  So a consistency judgement that
violates (P3) is reachable from skeleton inference in one step.

% =====================================================================
\section{The gap}\label{sec:gap}
% =====================================================================

Take the subject and counterpart
\[
  \syn{S} \;=\; \syn{\Scex}
  \qquad\qquad
  \bid{P} \;=\; \syn{\Pcex}
\]
under any $\Gamma$ --- in particular under $\Gamma = \cdot$.  A ghost
sits directly under a quantifier of its own subject, and the counterpart
uses that quantifier's variable.

\begin{factc}[the rules accept it]\label{fact:consis} \MC
  $\typconsis[\cdot]{\Gamma}{\syn{\Scex}}{\syn{\Pcex}}$.
\end{factc}

\begin{proof}
  \eqlab{ForallSyn} on $\alpha$, which extends the context to
  $\Gamma,\alpha$; then \eqlab{Ghost} with $\inh{P} \defeq \alpha$,
  whose only obligation is $\Gamma,\alpha\vdash\alpha$.
  (\coq{cap\_consis}.)
\end{proof}

\begin{factc}[the side condition rejects it]\label{fact:refine} \MC
  $\skctx[\cdot]{\Gamma}\Scex \fill \Pcex$ does \textbf{not} hold, for
  any $\Gamma$.
\end{factc}

\begin{proof}
  The only rule with a $\forall$ on both sides is \textsc{R-Forall}, and
  it moves $\alpha$ into $\Delta$.  The residual obligation is
  $\skctx[\alpha]{\Gamma,\alpha}\Hole\fill\alpha$, and the only rule
  with $\Hole$ on the left is \textsc{R-Ghost}, whose side condition
  reads $\{\alpha\}\cap\ftv{\alpha}=\emptyset$.
  (\coq{cap\_no\_refine}.)
\end{proof}

\noindent
So (P3) is a real restriction on the consistency \emph{rules}, not a
consequence of them.  The mechanism is the asymmetry (P1)/(P2) names:

\begin{center}\small
\begin{tabular}{@{}lll@{}}
  \toprule
  & where a $\forall$-bound $\alpha$ goes & may a ghost's fill use it? \\
  \midrule
  consistency, \eqlab{ForallSyn} & $\Gamma$ & yes --- \eqlab{Ghost} has no condition \\
  refinement, \textsc{R-Forall}  & $\Delta$ & no --- \textsc{R-Ghost} forbids $\Delta$ \\
  \bottomrule
\end{tabular}
\end{center}

\noindent
One caveat, because it is the exception that confirms the rule.  The
universal ghost $\TopHole$ behaves the \emph{other} way: its refinement
rule
$\skctx[\Delta]{\Gamma,\alpha}\TopHole.\inh{P}\fillsyn\bid{Q}
 \;\Longrightarrow\;
 \skctx[\Delta]{\Gamma}\TopHole.\inh{P}\fillsyn\syn{\forall\alpha.}\bid{Q}$
extends $\Gamma$, not $\Delta$.  A universal ghost \emph{is} allowed to
capture --- that is what generalisation means --- and any repair must
keep that.  Only $\forall$ (both colours) creates the conflict.

% =====================================================================
\section{Consequence 1: skeleton inference, as published, outruns the
  algorithm}\label{sec:inference}
% =====================================================================

Instantiate \patlab{Var} at $\Gamma = x : \Scex$.  The entry is
well-formed: a ghost under a quantifier is a legal skeleton.

\begin{factc}\label{fact:pat} \MC
  $\pati{x:\Scex}{x}{\syn{\Pcex}}$.
\end{factc}

\begin{proof}
  \patlab{Var} on $\Gamma\ni x:\Scex$ with \sublab{Infer} on
  \Cref{fact:consis}.  (\coq{cap\_d\_pat}.)
\end{proof}

\begin{factc}\label{fact:notreach} \MC
  Once the base rules carry (P3), the same judgement is
  \textbf{underivable}.
\end{factc}

\begin{proof}
  \patlab{Var} is the only rule for a variable, the singleton context
  pins the skeleton to $\Scex$, \sublab{Infer} is the only subtyping rule
  at $\M=\TI$, and its consistency premise now owes
  \Cref{fact:refine}.  (\coq{cap\_not\_reachable}.)
\end{proof}

\noindent
\Cref{fact:notreach} is not hypothetical: the mechanisation's
decorated judgement $\patp[\M]{\Theta}{M}{\bid{Q}}$ carries (P3) at every
consistency leaf of the embedded subtyping derivation, because that is
\textsc{Completeness and optimality of subtyping} needs (\Cref{sec:solutions}).  So the two systems
disagree, and \Cref{fact:pat} is on the wrong side of the disagreement:

\begin{factc}\label{fact:bridge} \MC
  There is no bridge
  $\patp[\M]{\Gamma}{M}{\bid{Q}} \Longrightarrow
   \patp[\M]{\Theta}{M}{\bid{Q}}$ from the published rules to the
  side-condition-carrying ones, even assuming $\Gamma$ well-formed and
  $M$ well-formed in $\Gamma$.  (\coq{d\_pat\_bridges\_false}.)
\end{factc}

\noindent
This is why the closed-program corollaries of \textsc{Completeness and
optimality of skeleton inference} and \textsc{\dots of type inference}
cannot at present be restated over the published declarative judgements: their completeness halves need exactly this
bridge.

% =====================================================================
\section{Consequence 2: \patlab{App}'s existential skeleton does not
  help}\label{sec:app}
% =====================================================================

\Cref{fact:bridge} refutes a \emph{pointwise} bridge, and one might hope
that is an artefact.  \patlab{App} and \sdtylab{App} do not fix the
argument's skeleton:
\[
\inferrule*[Lab=\sdtylab{App}]
{
  {\pati{\Gamma}{N}{\syn{P}}} \\
  {\Gamma \vdash  \inh{P} \fillsyn \bid{A}} \\
  \sdtypp[\FN \M]{\Gamma}{M}{\bid{A}\to \bid{B}} \\
  \sdtypc{\Gamma}{N}{\inh{A}}
}
{\sdtypp[\M]{\Gamma}{M\;N}{\bid{B}}}
\]
so a translation may answer with \emph{some other} $\syn{P'}$, provided
it still refines into the same $\bid{A}$.  It cannot.

\begin{factc}\label{fact:app} \MC
  At $\Gamma = x:\Scex$, $N = \Freeze{x}$ and
  $\bid{A} = \inh{\Pcex}$, the published rules give
  $\pati{\Gamma}{N}{\syn{\Pcex}}$ together with
  $\Gamma \vdash \inh{\Pcex}\fillsyn\inh{\Pcex}$; and there is
  \textbf{no} $\syn{P'}$ at all with
  $\pati{\Theta}{N}{\syn{P'}}$ and
  $\Theta \vdash \inh{P'}\fillsyn\inh{\Pcex}$.
  (\coq{d\_pat\_app\_package\_false}.)
\end{factc}

\begin{proof}[Proof sketch]
  \patlab{Freeze} is the only rule for $\Freeze{x}$ and it preserves the
  skeleton, so $\syn{P'}$ must come from \patlab{Var} at $x:\Scex$, whose
  side condition (P3) forces $\skctx[\cdot]{\Theta}\Scex\fill P'$.  By
  \textsc{R-Forall} and \textsc{R-Ghost} that makes $P' = \forall\alpha.Q$
  with $\alpha\notin\ftv{Q}$ and $Q$ scoped in $\Theta$.  The refinement
  premise of \sdtylab{App} demands
  $\skctx[\cdot]{\Theta}\inh{\forall\alpha.Q}\fillsyn\inh{\Pcex}$, whose
  body obligation is $\skctx[\alpha]{\Theta}\inh{Q}\fillsyn\alpha$ ---
  and every rule that can produce a bare variable on the right
  (\textsc{R-Var}, \textsc{R-Ghost}, and the $\TopHole$-removal chain
  above them) puts $\alpha$ in $\Theta$.  Both cannot hold.
\end{proof}

\noindent
The two demands read the \emph{same} binder in opposite directions: the
side condition wants $P'$'s body free of $\alpha$, the App refinement
wants it to \emph{be} $\alpha$.  With $\bid{A}$ free there is an answer
($P' = \Scex$ against $\bid{A} = \inh{\forall\alpha.\tunit}$), so the
refutation genuinely uses the fact that \sdtylab{App}'s $\bid{A}$ is
pinned by the function premise.

% =====================================================================
\section{Why the side condition is at the leaves, and what the solution
  notions do about it}\label{sec:solutions}
% =====================================================================

The paper's three problem/solution definitions
(\coq{algorithmic.tex:392}, \coq{:559}, \coq{:1423}) all attach a
refinement datum \emph{at the root}.  For subtyping:

\begin{quote}\itshape
  A subtyping problem is a tuple $(\Theta_0;\N;\M;P_0;P)$ \dots\ A
  solution to it is a tuple $(\inc[\theta_1]{\Theta_0}{\Theta_1};P_1)$
  such that $\typsub[\N][\M][]{\Theta_1}{\syn{\theta_1 P}}{\bid{P_1}}$
  where $\Theta_1 \vdash \inh{\theta_1 P_0}\fillsyn\bid{P_1}$.
\end{quote}

\noindent
Two facts about that root datum, both already recorded in the
mechanisation:

\begin{enumerate}[leftmargin=*]
\item \textbf{It does not imply (P3).}  Take the configuration of
  \Cref{sec:gap} with $P_0 = \Hole$, $\N=\SK$, $\M=\TI$, $\theta_1 =
  \mathit{id}$.  The join holds by \sublab{Infer} and
  \Cref{fact:consis}; the root datum
  $\skctx[\cdot]{\Theta_1}\Hole\fill\Pcex$ holds by \textsc{R-Ghost},
  since $\Delta = \cdot$ there and the condition
  $\cdot\cap\ftv{\Pcex}=\emptyset$ is vacuous.  So both solution
  conditions pass while (P3) fails at the subject.  \ARG

\item \textbf{Moving (P3) to the root instead would be wrong in the other
  direction.}  Refinement is a lockstep relation, so a root-level
  subject-refinement premise excludes every solution whose two sides
  disagree at a head --- that is, precisely \sublab{ForallL},
  \sublab{ForallR} and \sublab{UnivGhostR}, the rules instantiation and
  generalisation live in.  \MC
  (\coq{cex/cex\_p7\_univ\_ghost\_r\_vacuous.v},
  \coq{cex/cex\_p7\_third\_premise\_forall.v}.)
\end{enumerate}

\noindent
Hence the mechanisation carries (P3) where consistency actually
occurs --- at the \sublab{Infer}, \sublab{Check} and
\sublab{UnivGhostR} \emph{leaves} of a subtyping derivation, and at the
domain premise of \sublab{Arrow}, which is itself a consistency call.
That placement is not a strengthening of the paper: it is the paper's own
sentence (P3), applied to each consistency judgement the derivation
contains.

\medskip\noindent
For the record, the same reading is what makes the augmented
\textsc{Completeness of consistency} (\coq{consistency-proofs.tex:2230--2242})
correct: its third solution condition \emph{is} (P3), and the published
two-condition form is false.

% =====================================================================
\section{What survives the repair, and what it is for}\label{sec:survives}
% =====================================================================

Restoring (P3) to the rules makes \Cref{fact:consis} underivable, and
therefore \Cref{fact:pat} too.  It is worth being explicit about what
that does to the rest, because the two halves of this note have
different standing.

\begin{center}\small
\begin{tabular}{@{}p{7.2cm}p{7.2cm}@{}}
  \toprule
  \textbf{Survives the repair} & \textbf{Retracted by the repair} \\
  \midrule
  \Cref{fact:consis} + \Cref{fact:refine}: the consistency \emph{rules}
  derive a pair that refinement rejects.  A statement about the published
  rules; it is the proof that (P3) is a genuine restriction rather than
  an derivable convenience, and it is why (P3) may not be dropped.
  &
  \Cref{fact:pat}, \Cref{fact:bridge}, \Cref{fact:app}: statements about
  a skeleton-inference judgement read \emph{without} (P3).  After the
  repair there is no such judgement, and all three become historical
  witnesses --- the record of why the presupposition could not be
  deferred. \\
  \bottomrule
\end{tabular}
\end{center}

\noindent
This is the same shape as the earlier \soleqlab{FlexR} note in this
directory (\coq{docs/refutations.tex}), whose two refutations were
retracted by the repair they motivated and kept as its justification.

\medskip\noindent
Two further points on standing.  First, \emph{no published statement is
being contradicted}: (P3) is the paper's own sentence, so the repair is a
transcription, and the facts that retract are facts about the
mechanisation's debt, not about \Calc.  Second, the repair is confined to
the \emph{skeleton} sort.  Type consistency
$\typconsis[\Psi]{\Gamma}{\syn{T}}{\bid{A}}$ ranges over types, which
have no ghosts, so (P3) holds there for free --- which is exactly why the
type-sorted half of the bridge is already closed and the skeleton-sorted
half is not.  \sdtylab{Unit}, \sdtylab{Var} and \sdtylab{Anno} are
therefore untouched; only \sdtylab{App}, through its
$\pati{\Gamma}{N}{\syn{P}}$ premise, ever reaches the affected judgement.

% =====================================================================
\section{The repair}\label{sec:repair}
% =====================================================================

Put (P3) on the rules.  Two renderings; they differ in where the work
goes, not in what they exclude.

\paragraph{(R-a) Carry it as a premise.}
Add $\skctx[\cdot]{\Gamma}\erase{S}\fill P$ to
$\typconsis[\Psi]{\Gamma}{\syn{S}}{\bid{P}}$ --- as the paper's sentence
literally says --- and, in the mechanisation, to the consistency-bearing
rows of skeleton subtyping, mirroring what the decorated judgement
already does.  This is a transcription of (P3) and nothing more, so it
cannot exclude anything the paper does not already exclude.  Cost: a
premise on the rows, and one new obligation wherever a consistency
derivation is \emph{built}.

\paragraph{(R-b) Make it derivable, by giving consistency the $\Delta$
that refinement has.}
Index skeleton consistency by $\Delta$,
$\typconsis[\Psi]{\skctx[\Delta]{\Gamma}}{\syn{S}}{\bid{P}}$, and:
\begin{itemize}[leftmargin=*,itemsep=.2ex]
  \item the four $\forall$ rows (\eqlab{ForallInh}, \eqlab{ForallSyn},
    tagged and untagged) put $\alpha$ into $\Delta$, not $\Gamma$;
  \item \eqlab{Var} reads $\Gamma,\Delta\ni\alpha$, as
    \textsc{R-Var} and skeleton well-formedness already do;
  \item \eqlab{Ghost}'s $\inh{P}$ is well-formed with $\Delta$
    \emph{cleared}, matching \textsc{R-Ghost};
  \item \eqlab{UnivGhost1} keeps extending $\Gamma$ --- see the caveat
    at the end of \Cref{sec:gap}.
\end{itemize}
Then (P3) is a \emph{theorem} about consistency rather than an
assumption, at arbitrary subjects, and the leaf premises of the
decorated judgement can be deleted rather than mirrored.

\paragraph{Which.}
(R-a) is the faithful minimum and is what this note recommends stating in
the paper; (R-b) is the better mechanisation, because it removes a
premise instead of adding one and it makes the two judgements' scoping
disciplines visibly the same.  They should agree, but that has to be
checked, and there is one known place where they might not: the
mechanised \textsc{R-TopHoleKeep} carries an explicit
$\guarded{\bid{Q}}$ that the paper's rule has commented out
(\coq{declarative.tex:1508}, \coq{docs/WF-SIDE-CONDITIONS.md} \S3), so
(R-b)'s derivation of (P3) at $\eqlab{UnivGhostG}$ and $\eqlab{UnivGhost}$
needs that premise or the paper's version of the rule.

\paragraph{What the repair buys.}
\Cref{fact:pat} stops being derivable, so \Cref{fact:bridge,fact:app}
evaporate; the published and decorated skeleton-inference judgements
coincide; and the completeness halves of the two inference theorems can be
restated over the published declarative system, which is what the
closed-program corollaries want.

\paragraph{What the repair costs, expressiveness-wise.}
Exactly the derivations \Cref{fact:refine} rejects: a ghost filled with a
variable bound by a quantifier of its own skeleton.  No algorithmic run
produces such a fill --- algorithmic consistency solves ghosts with
flexible variables scoped before the rigid is introduced --- so nothing
the implementation accepts is lost.  Whether any \emph{source program}
relies on one is the authors' question to answer; this note has not found
one, and \Cref{sec:app}'s analysis suggests the configuration is hard to
reach from a well-typed term, since a $\forall$-bound ghost can only be
\emph{refined into} a skeleton, never created by one.

% =====================================================================
\section{Status}\label{sec:status}
% =====================================================================

\noindent
Against \coq{mechanization/} at commit \coq{2cb0179} (branch
\coq{remech}): \coq{make coq} exit 0, no \texttt{Admitted}, no
\texttt{Axiom}, and every \texttt{Print Assumptions} in \coq{README.v}
reads \emph{closed under the global context}.

\medskip\noindent
Nothing in \Cref{sec:gap,sec:inference,sec:app} depends on a
mechanisation-only definition: \Cref{fact:consis,fact:refine} are about
the paper's own consistency and refinement rules, and
\Cref{fact:pat,fact:notreach,fact:bridge,fact:app} about the paper's
skeleton inference and \sdtylab{App}.  The one place a
mechanisation-internal judgement appears is the
side-condition-carrying skeleton inference of
\Cref{fact:notreach,fact:app}, and \Cref{sec:solutions} argues that
judgement \emph{is} the paper's, read with (P3) in force.

\end{document}
